% !PRINTED 249
\weberpage{249}
If one counts zero among the squares, one obtains the theorem:

\gesperrt{6. If $p = 2$, all numbers in $\mathfrak{C}$ are squares. If $p$ is odd, there are $\frac{1}{2}(p^n + 1)$ squares, $\frac{1}{2}(p^n - 1)$ non-squares in $\mathfrak{C}$.}

For odd $p$ a primitive root must always be a non-square, and in the sequence (8) the numbers with even exponents are the squares, those with odd exponents the non-squares.

The product and the quotient of two squares are each evidently again squares.

From this it follows that the product of a non-square and a square different from zero is a non-square. For if the product $\alpha\beta$ is a square and the one factor $\alpha$ is a square, then $\beta = \alpha\beta : \alpha$ must also be a square.

From this one concludes further that the product of two non-squares is a square. For if $\beta$ denotes any non-square, let the factor $\xi$ in the product $\beta\xi$ run through all the numbers of $\mathfrak{C}$ different from zero. The product then runs through the same sequence of numbers, and since $\beta\xi$ is a non-square for every square $\xi$, it must be a square for every non-square $\xi$.

All this also follows very simply from the representation (8) of the numbers of $\mathfrak{C}$ by a primitive root.

We close these general considerations with the theorem, which we shall need later:

\gesperrt{7. Every non-square number is the sum of two squares.}

Here $p$ is to be assumed odd, since only then are there non-square numbers.

First of all, every number $\beta$ can then be represented as the difference of two squares, as the identity

\[
\left(\beta + \tfrac{1}{4}\right)^2 - \left(\beta - \tfrac{1}{4}\right)^2 = \beta
\]

shows. Now if $-1$ is a square in $\mathfrak{C}$, one sets

\[
\left(\beta + \tfrac{1}{4}\right)^2 = \xi^2, \quad -\left(\beta - \tfrac{1}{4}\right)^2 = \eta^2,
\]

and obtains

\[
\beta = \xi^2 + \eta^2.
\]

But if $-1$, and therefore also $p - 1$, is a non-square, then consider the sequence of the natural numbers $1, 2, \ldots, p - 1$, whose
%CONT>>
